\section{The object, the target, and the exact record}

This part fixes the object of Erd\H{o}s Problem \#249, states exactly what is
open about it, and records every unconditional numeral, identity and
Lean-checked bound the corpus currently holds. Nothing in this part decides
irrationality. Every result below is either an identity (no hypothesis, holds
unconditionally), a finite computation (holds up to an explicitly stated
bound), or is flagged \ev{Cited}/\ev{Open} where that is the honest status.

\subsection{The object}

\begin{defn}[The Erd\H{o}s--249 constant]
\label{defn:S}
Let $\varphi$ denote Euler's totient function. Define
\[
S \;:=\; \sum_{n \ge 1} \frac{\varphi(n)}{2^{n}} \;=\; \sum_{n : \mathbb{N}} \frac{\varphi(n)}{2^{n}},
\]
the two forms coinciding because $\varphi(0)=0$; the second, $\mathbb{N}$-indexed
form is the one carried in the Lean source. The series converges absolutely
since $\varphi(n) = O(n)$. \emph{Erd\H{o}s \#249 asks whether $S$ is irrational.
This is OPEN.} Nothing in this paper decides it; every claim below is either an
unconditional identity, a finite computation with an explicitly stated range,
or is marked \ev{Cited}/\ev{Open}.
\end{defn}

\begin{rem}[Status]
No proof or disproof of $\mathrm{Irrational}(S)$ exists anywhere in the corpus,
formal or informal. What exists is: one large unconditional finite denominator
exclusion (\S\ref{ssec:farey}), several exact reformulations of $S$ that
relocate the same open question onto different coordinates without touching
its truth value (\S\ref{ssec:coprime}--\S\ref{ssec:lambert}), and a certificate
apparatus whose cofinal supply obligation is the precise open target
(recorded in full in Part~1 of this paper; only its finite, checked instances
are catalogued here, \S\ref{ssec:certtable}).
\end{rem}

\subsection{The binary-digit reduction}

\begin{prop}[Digit-shift identity]
\label{prop:shift}
For every $N : \mathbb{N}$,
\[
2^{N} \cdot S \;=\; \Phi_N + R_N,
\]
where $\Phi_N := \sum_{n \le N} \varphi(n) \cdot 2^{N-n} \in \mathbb{N}$ (the
integer prefix) and $R_N := \sum_{j \ge 0} \varphi(N+1+j)/2^{j+1}$ (the
fractional tail, \texttt{totientTail} in Lean).
\coord{binary-digit} \scale{uniform} \ev{Lean}
\lean{two_pow_mul_totient_series_eq}{TotientTailPeriodKiller.lean:150}
\end{prop}

\begin{proof}[Consequence]
Rationality and eventual periodicity of the base-2 expansion of $S$ is
\emph{exactly} a statement about the tail $R_N$: $S \in \mathbb{Q}$ forces
(by Euler's theorem applied to the odd part of the denominator) a period
$h>0$ and pre-period $N_0$ with $R_{N+h}-R_N \in \mathbb{Z}$ for all
$N \ge N_0$ (\texttt{eventual_period_of_not_irrational},
\lean{eventual_period_of_not_irrational}{TotientTailPeriodKiller.lean:358},
\ev{Lean}, explicit witness $h=\varphi(\mathrm{oddPart}(r.\mathrm{den}))$,
$N_0=v_2(r.\mathrm{den})$). Nothing in the corpus supplies the converse
cofinally; this is the certificate wall documented in full in Part~1 and is
not restated here beyond this pointer.
\end{proof}

\subsection{The unconditional denominator floor}
\label{ssec:farey}

The strongest \emph{unconditional} fact the corpus holds about $S$ is a lower
bound on its denominator should it happen to be rational. It is obtained by a
completely elementary Farey/mediant argument, entirely free of any certificate
or period apparatus, and is therefore logically independent of the open
certificate-supply obligation.

\begin{lem}[Farey gap, fully general]
\label{lem:farey}
For integers $a,b,c,d,r,s$ with $b>0$, $d>0$, $bc-ad=1$ (i.e.\ $a/b$ and
$c/d$ are unimodular Farey neighbours), and $as < rb$, $rd < cs$ (i.e.\ $r/s$
lies strictly between them): $b+d \le s$.
\coord{farey} \scale{n/a} \ev{Lean}
\lean{farey_gap}{GapFareyBound.lean:51}
\end{lem}

This lemma is a classical Stern--Brocot fact with zero totient or Mersenne
content; it is the generic engine underneath every Farey-gap bound in the
corpus, for either open problem.

\begin{prop}[The wave-17 gap certificate at window $K=240$]
\label{prop:gapwindow}
Let $V$ be the explicit committed totient residue for window $(N,K)=(1,240)$.
For every $q : \mathbb{N}$ with
\[
0 < q \;\le\; Q_0 := 79\,639\,646\,646\,701\,375\,323\,355\,774\,875\,831\,053 \;\;(\approx 7.96\times 10^{34}),
\]
\[
(q \cdot V) \bmod 2^{240} \;+\; 243\,q \;<\; 2^{240}.
\]
This bound is \emph{sharp}: $q = Q_0+1 = 79\,639\,646\,646\,701\,375\,323\,355\,774\,875\,831\,054$
is the exact first failing denominator, exhibited as the mediant of two
explicit unimodular Farey neighbours.
\coord{farey} \scale{bounded} \ev{Lean}
\lean{gap_check_window_1_240_le_79639646646701375323355774875831053}{GapFareyBound.lean:176}
\lean{gap_check_window_1_240_first_failure}{GapFareyBound.lean:225}
\end{prop}

\begin{thm}[Erd\H{o}s \#249 denominator exclusion --- the headline unconditional result]
\label{thm:denom-record}
For every $p \in \mathbb{Q}$ with reduced denominator $p.\mathrm{den} \le Q_0$,
\[
S \;\neq\; p.
\]
Equivalently: \emph{if $S$ is rational, its reduced denominator exceeds}
$Q_0 \approx 7.96 \times 10^{34}$.
\coord{farey} \scale{bounded} \ev{Lean}
\lean{tsum_totient_div_pow_two_ne_ratCast_of_den_le_79639646646701375323355774875831053}{CertificateKernel.lean:18572}
\end{thm}

\begin{rem}[What this does and does not say]
Theorem~\ref{thm:denom-record} is a complete, unconditional finite fact: no rational
of small denominator equals $S$. It says nothing about arbitrarily large
denominators and is not itself a route to irrationality. The docstring behind
Proposition~\ref{prop:gapwindow} states the honest open question this leaves:
whether $\sup_K (b+d)(K)$ (the growing analogue of $Q_0$ as the window $K$
grows) is unbounded as $K \to \infty$ --- which would close \#249 through this
theorem's consumer, entirely independently of the certificate-supply
obligation of Proposition~\ref{prop:shift}. \ev{Open}; not claimed or proved
anywhere in the corpus.
\end{rem}

\begin{rem}[Ladder of prior rungs]
$Q_0$ is the current end of an explicit sequence of increasingly wide Farey
windows computed by the same mechanism: $4838 \to 2^{22} \to 2.49\times10^{17}$
(window $K=120$) $\to Q_0$ (window $K=240$). Each rung is a finite,
independently checked instance of Lemma~\ref{lem:farey} and
Proposition~\ref{prop:gapwindow}'s pattern at a larger $K$; none is claimed to
extrapolate.
\end{rem}

\subsection{The same record, transported: the M\"obius-square and coprimality-probability forms}
\label{ssec:coprime}

The exact finite Farey record behind Theorem~\ref{thm:denom-record} is not tied to
the $\varphi(n)/2^n$ presentation of $S$; it transfers verbatim to two other
exact reformulations of the same constant, at exactly half the bound.

\begin{prop}[Fair-coin coprimality form]
\label{prop:coprime}
Let $X,Y$ be independent random variables with $\Pr(X=n)=\Pr(Y=n)=2^{-n}$
for $n \ge 1$ (independent fair-coin waiting times). Then
\[
S \;=\; \tfrac12 \;+\; \Pr\bigl(\gcd(X,Y)=1\bigr)
\;=\; \tfrac12 \;+\; \sum_{\substack{a,b\ge 1\\ \gcd(a,b)=1}} 2^{-(a+b)}.
\]
Equivalently, on the visible lattice: summing $2^{-(a+b)}$ over the half-open
coprime pairs ($a\ge1$, $b\ge0$, $\gcd(a,b)=1$) recovers $\sum_n \varphi(n)/2^n$
exactly, with no boundary correction, because the visible-point count on the
half-open antidiagonal at height $n$ equals $\varphi(n)$ for every $n$,
including $n=0,1$.
\coord{probability} \scale{n/a} \ev{Lean}
\lean{tsum_visible_coprime_pairs_eq_totient_series}{CertificateKernel.lean:18732}
\lean{totient_series_eq_half_add_visible_coprime_pairs}{CertificateKernel.lean:18745}
\end{prop}

\begin{prop}[The gcd-layer normalisation]
\label{prop:gcdlayer}
For independent fair-coin waiting times as above, $\sum_{g\ge1}\Pr(\gcd(X,Y)=g)=1$
exactly; and for every $d>0$, $\Pr(d\mid X \wedge d\mid Y) = 1/(2^d-1)^2$.
\coord{probability} \scale{uniform} \ev{Lean}
\lean{tsum_pos_coprime_inv_mersenne_eq_one}{GcdMomentCalculus.lean:349}
\lean{tsum_pos_pair_both_dvd_half_eq_inv_mersenne_sq}{GcdMomentCalculus.lean:266}
\end{prop}

\begin{thm}[Denominator exclusion for the coprimality-probability form]
\label{thm:denomcoprime}
Let
\[
Q_1 := \left\lfloor \frac{Q_0}{2} \right\rfloor = 39\,819\,823\,323\,350\,687\,661\,677\,887\,437\,915\,526.
\]
For every $a\in\mathbb{Z}$, $d\in\mathbb{N}$ with $0<d\le Q_1$: the visible
coprime-pair probability $\Pr(\gcd(X,Y)=1)$ is not equal to $a/d$.
\coord{farey} \scale{bounded} \ev{Lean}
\lean{tsum_visible_coprime_pairs_ne_int_div_of_den_le_39819823323350687661677887437915526}{CertificateKernel.lean:18760}
\end{thm}

\begin{thm}[Denominator exclusion for the M\"obius-square form]
\label{thm:denommobsq}
With $Q_1$ as in Theorem~\ref{thm:denomcoprime}: for every $a\in\mathbb{Z}$,
$d\in\mathbb{N}$ with $0<d\le Q_1$, the signed series
$T := \sum_{d\ge1} \mu(d)/(2^d-1)^2 = S - \tfrac12$ (see \S\ref{ssec:mobius})
is not equal to $a/d$.
\coord{farey} \scale{bounded} \ev{Lean}
\lean{tsum_moebius_div_two_pow_sub_one_sq_ne_int_div_of_den_le_39819823323350687661677887437915526}{CertificateKernel.lean:18675}
\end{thm}

\begin{rem}[Why the bound halves, and why this is not new information]
$Q_1 = \lfloor Q_0/2 \rfloor$ arithmetically: $Q_0$ is odd
($Q_0 = 79\,639\,646\,646\,701\,375\,323\,355\,774\,875\,831\,053$), so
$Q_0/2 = 39\,819\,823\,323\,350\,687\,661\,677\,887\,437\,915\,526.5$ and
$Q_1$ is its floor. The halving is the cost of transporting the known bound
through the affine shift
$S = \tfrac12 + T$ (resp.\ $S = \tfrac12 + \Pr(\gcd(X,Y)=1)$) on a
denominator-exclusion statement: excluding all denominators $\le Q_1$ for $T$
follows from excluding all denominators $\le Q_0$ for $S$ (a denominator-$d$
value of $T$ with $d\le Q_1$ yields a
denominator dividing $2d\le Q_0$ for $S$). Theorems~\ref{thm:denomcoprime}
and~\ref{thm:denommobsq} are therefore the \emph{same} finite Farey record as
Theorem~\ref{thm:denom-record}, transported through Proposition~\ref{prop:coprime}
and Proposition~\ref{prop:mobsq} respectively --- not independent evidence.
The converse finite implication is not asserted: subtracting $1/2$ can double
a denominator, so the $Q_1$ exclusion for $T$ alone need not recover the full
$Q_0$ exclusion for $S$.
\end{rem}

\subsection{The M\"obius--Mersenne identity and why bounded coefficients matter}
\label{ssec:mobius}

\begin{prop}[M\"obius-square reduction]
\label{prop:mobsq}
\[
S \;=\; \sum_{n\ge1}\frac{\varphi(n)}{2^n} \;=\; \frac12 \;+\; \sum_{d\ge1}\frac{\mu(d)}{(2^d-1)^2},
\]
where $\mu$ is the M\"obius function, so $\mu(d)\in\{-1,0,1\}$ for every $d$.
Consequently \emph{Erd\H{o}s \#249 $\iff$ $T:=\sum_{d\ge1}\mu(d)/(2^d-1)^2 \notin \mathbb{Q}$},
the single reduced target every other result in this subsection and
\S\ref{ssec:lambert} feeds.
\coord{mobius-mersenne} \scale{n/a} \ev{Lean}
\lean{totientSeries_eq_pnat_half_pow}{SquaredMersenneDiagonalEnclosure.lean:72}
\lean{tsum_totient_half_pow_eq_half_add_moebius_sq}{MersenneLambertLadder.lean:665}
\end{prop}

\begin{rem}[Why the bounded coefficients matter --- the Erd\H{o}s-1948 regime]
The identity of Proposition~\ref{prop:mobsq} rewrites $S$ (equivalently $T$)
as a \emph{M\"obius-twisted Lambert-squared series}: the numerator weight
$\mu(d)$ is bounded, $|\mu(d)|\le1$ for every $d$, uniformly in $d$. This
places $T$ in exactly the coefficient regime of the classical Erd\H{o}s (1948)
near-integer irrationality criterion and of the level-1 sibling identity
$L(\mu):=\sum_d \mu(d)/(2^d-1) = \tfrac12$ (rational, trivially) alongside
$L(1) = \sum_d 1/(2^d-1) = E$, the Erd\H{o}s--Borwein constant, which
\emph{is} proved irrational in this same kernel
(\lean{irrational_erdosBorwein_series}{CertificateKernel.lean:8523}, \ev{Lean}).
This is the opposite regime from Proposition~\ref{prop:shift}'s
\textbf{binary-digit coordinate}, where the corresponding weight satisfies
$0\le\varphi(n)\le n$ and is unbounded.  The function $\varphi$ has average
order $6n/\pi^2$, equivalently
$\sum_{k\le x}\varphi(k)\sim 3x^2/\pi^2$, but there is no pointwise estimate
$\varphi(n)=\Theta(n)$. A
near-integer/Dirichlet-approximation argument of Erd\H{o}s-1948 shape
(formalised generically as
\lean{irrational_of_int_mul_near_int}{CertificateKernel.lean:6308} and its
base-power specialisation
\lean{irrational_of_pow_mul_near_int}{CertificateKernel.lean:6337}, both
\ev{Lean}, \coord{n/a}, fully coordinate-free) has a genuine chance of
transferring to $T$ precisely because its weight is bounded, in a way it does
not have a chance of transferring directly to the raw $\varphi(n)/2^n$ series.
No such transfer is proved; \S\ref{ssec:mobius} below (cross-referenced
here, developed in Part~2 of this paper) records exactly why the transfer has
so far failed (the ``$\mu$-pollution'' obstruction) rather than merely
asserting the analogy.
\end{rem}

\subsection{The squared-Lambert gcd-moment identities}
\label{ssec:lambert}

\begin{prop}[Squared-Lambert transfer engine]
\label{prop:lambertengine}
For $w:\mathbb{N}\to\mathbb{R}$ with $|w(d)|\le d$ for all $d>0$, and $0\le r<1$:
\[
\sum_{d\ge1} w(d)\left(\frac{r^d}{1-r^d}\right)^2 \;=\; \sum_{n\ge1}\left(\sum_{e\mid n} w(e)\Bigl(\tfrac{n}{e}-1\Bigr)\right) r^n.
\]
\coord{mobius-mersenne} \scale{uniform} \ev{Lean}
\lean{tsum_lambert_linear_weight_sq_pure}{GcdMomentCalculus.lean:105}
\end{prop}

This one identity, at $r=1/2$, specialises to every squared-Lambert rung the
corpus computes; two instances are exact and directly relevant to \#249's
weight structure:

\begin{prop}[The known $\zeta_q$-rung]
\label{prop:zetaq}
\[
\sum_{d\ge1} \frac{1}{(2^d-1)^2} \;=\; \sum_{n\ge1} \frac{\sigma(n)-\tau(n)}{2^n} \;=\; \zeta_q(2)-\zeta_q(1) \text{ at } q=\tfrac12,
\]
where $\sigma$ is the sum-of-divisors function and $\tau$ the number-of-divisors
function. The \emph{identity} is machine-checked (\ev{Lean}); irrationality of
the \emph{value} $\zeta_q(2)-\zeta_q(1)$ is \ev{Cited} (Postelmans--Van Assche
$q$-Pad\'e), \emph{not} formalised in this corpus.
\coord{mobius-mersenne} \scale{n/a}
\lean{tsum_one_div_mersenne_sq_eq_sigma_sub_tau_series}{GcdMomentCalculus.lean:216}
\end{prop}

\begin{prop}[The Pillai/gcd-moment rung]
\label{prop:pillai}
\[
\sum_{d\ge1} \frac{\varphi(d)}{(2^d-1)^2} \;=\; \sum_{n\ge1} \bigl(P(n)-n\bigr)\cdot 2^{-n} \;=\; \mathbb{E}[\gcd(X,Y)],
\]
where $P(n) := \sum_{e\mid n}\varphi(e)\cdot(n/e) = (\varphi * \mathrm{Id})(n)$
is Pillai's gcd-sum function and $X,Y$ are the independent fair-coin waiting
times of Proposition~\ref{prop:coprime}. The identity itself is
machine-checked (\ev{Lean}); the value $\mathbb{E}[\gcd(X,Y)]$ is \ev{Open} ---
a cousin rung to \#249, not \#249 itself.
\coord{mobius-mersenne} \scale{n/a}
\lean{tsum_totient_div_mersenne_sq_eq_gcd_moment_series}{GcdMomentCalculus.lean:235}
\end{prop}

\begin{rem}[The level mirror]
Writing $L(f):=\sum_d f(d)/(2^d-1)$ (level 1) and $L_2(f):=\sum_d f(d)/(2^d-1)^2$
(level 2): $L(\mu)=\tfrac12$ (rational, trivial), $L(1)=E$ (Erd\H{o}s--Borwein,
irrational, \ev{Lean}), $L(\varphi)=2$ (rational); $L_2(\mu)=S-\tfrac12$
(\ev{Open}, this is \#249), $L_2(1)=\zeta_q(2)-\zeta_q(1)$ (\ev{Cited}
irrational, Proposition~\ref{prop:zetaq}), $L_2(\varphi)=\mathbb{E}[\gcd(X,Y)]$
(\ev{Open}, Proposition~\ref{prop:pillai}). At level 1 the M\"obius rung is
trivial and the $\zeta$-rung is the hard classical case; at level 2 the
$\zeta$-rung is known and the M\"obius rung \emph{is} \#249. M\"obius
projection is the one wall repeated at both levels of the ladder.
\end{rem}

\subsection{The exact certificate record: the contiguous band through $t\le82$}
\label{ssec:certtable}

The certificate apparatus itself (the predicate \texttt{certifiedKill}, its
soundness/completeness theorems, and the cofinal supply obligation that is the
actual open target) is developed in full in Part~1 of this paper. What is
recorded here is a representative historical set of machine-checked anchors,
together with the current aggregate theorem closing every scale $t\le82$.
The declaration inventory, rather than this table, is authoritative for all
individual certificate shards; no extrapolation beyond the stated band is
implied.

\paragraph{Certificate object.} For $h,N,L : \mathbb{N}$, define the window
discrepancy $\Delta_{h,N,L} := \sum_{j<L} (\varphi(N+h+1+j)-\varphi(N+1+j))\cdot 2^{L-1-j} \in \mathbb{Z}$
(\lean{windowDiscrepancy}{TotientTailPeriodKiller.lean:66}, \ev{Lean}), and
\[
\mathrm{Sep}(h,N,L) \;:\equiv\; (N+h+L+2) < \Delta_{h,N,L} \bmod 2^{L} < 2^{L}-(N+h+L+2),
\]
(\lean{certifiedKill}{TotientTailPeriodKiller.lean:72}, \ev{Lean}, decidable). By
completeness (\lean{exists_certifiedKill_iff_tail_diff_notMem_int}{LcmConeFlatness.lean:316},
\ev{Lean}), $(\exists L,\ \mathrm{Sep}(h,N,L)) \iff R_{N+h}-R_N \notin \mathbb{Z}$
for every $h,N$: an unbounded (cofinal, over $h$ and $N$) supply of
$\mathrm{Sep}$ is exactly equivalent to $\mathrm{Irrational}(S)$. The table
below records the principal $(h,N,L)$-shaped and $t$-shaped anchors; none of
them is cofinal, and none is presented as deciding \#249.

\small
\begin{longtable}{p{0.20\linewidth} p{0.40\linewidth} p{0.13\linewidth} p{0.20\linewidth}}
\textbf{Anchor} & \textbf{Exact statement of what was checked} & \textbf{Scale} & \textbf{Site} \\
\hline
\endhead
Fixed-window deposit, depth 16
&
$\mathrm{Sep}(h,12,16)$ holds for every $h \in \{1,\dots,8\}$ (by \texttt{decide},
$256$ totient values below $37$); consequently $S \ne r$ for every
$r\in\mathbb{Q}$ with $1\le h\le 8$ and $r.\mathrm{den} \mid 2^{12}(2^h-1)$.
&
\scale{fixed}
&
\lean{certifiedKill_all_small}{TotientTailPeriodKiller.lean:404}
\lean{totient_series_ne_rat_of_den_dvd}{TotientTailPeriodKiller.lean:416}
\\
Fixed-window deposit, depth 9, wider net
&
$\mathrm{Sep}(h,14,9)$ holds for every
$h\in\{1,\dots,16\}$: $S \ne r$ for every $r\in\mathbb{Q}$ with $1\le h\le16$
and $r.\mathrm{den}\mid 2^{14}(2^h-1)$.
&
\scale{fixed}
&
\lean{totient_series_ne_rat_of_den_dvd_pow_two_mul_mersenne_upto_sixteen}{CertificateKernel.lean:19078}
\\
Diagonal-pincer certificates, base 12 scales
&
$\mathrm{Sep}\bigl(H_t, H_t, D(t)\bigr)$, $H_t := \mathrm{lcm}(1,\dots,t)$,
holds (by explicit \texttt{decide}/\texttt{norm_num} on checked
\texttt{Nat.totient} values, via factored prime-power blocks with
Lucas-primality certificates) for $t \in \{1,2,3,4,5,7,8,9,11,13,16,17\}$
with certificate depths $D(t) \in \{6,5,7,7,9,14,15,14,21,22,23,26\}$
respectively (in the same order).
&
\scale{fixed}
&
\lean{certifiedKill_diagonal_all_imported}{DiagonalPincerCertificates.lean:2860}
\lean{diagonalPincerCertificateScales}{DiagonalPincerCertificates.lean:2842}
\lean{diagonalPincerKillDepth}{DiagonalPincerCertificates.lean:2844}
\\
Diagonal-pincer certificates, extended scales through $t=64$
&
The same predicate $\mathrm{Sep}(H_t,H_t,D(t))$ is additionally checked, one
sibling module per scale, at
$t \in \{19,\allowbreak23,\allowbreak25,\allowbreak27,\allowbreak29,\allowbreak31,
\allowbreak32,\allowbreak37,\allowbreak41,\allowbreak43,\allowbreak47,\allowbreak49,
\allowbreak53,\allowbreak59,\allowbreak61,\allowbreak64\}$
(16 further explicit values). Together with the base 12 scales above this is
\textbf{28 explicit values of $t$ in total}, the complete list being
$\{1,\allowbreak2,\allowbreak3,\allowbreak4,\allowbreak5,\allowbreak7,
\allowbreak8,\allowbreak9,\allowbreak11,\allowbreak13,\allowbreak16,\allowbreak17,
\allowbreak19,\allowbreak23,\allowbreak25,\allowbreak27,\allowbreak29,\allowbreak31,
\allowbreak32,\allowbreak37,\allowbreak41,\allowbreak43,\allowbreak47,\allowbreak49,
\allowbreak53,\allowbreak59,\allowbreak61,\allowbreak64\}$.
This does \emph{not} establish $\mathrm{Sep}(H_t,H_t,\cdot)$ for infinitely
many $t$; these 28 deposits are a historical strict subset of the contiguous
band in the next row.
&
\scale{fixed}
&
Per-scale Lean modules; endpoint
\lean{certifiedKill_diagonal_t64}{DiagonalPincerCertificateT64Endpoint.lean:1928}
\\
Contiguous lcm-diagonal band through $t\le82$
&
For every natural $t\le82$ there exists a depth $L$ with
$\mathrm{Sep}(H_t,H_t,L)$.  This closes every scale in the finite interval,
including plateau transfers, with no holes.  The next lcm jump is at the prime
$83$; no certificate at $t=83$ is claimed, and any such certificate must have
depth at least $125$.
&
\scale{bounded}
&
\lean{ErdosProblems.Skip.LadderT67.exists_diagonalKill_le_82}{ErdosProblems/Skip/LadderT67.lean:71264}
\lean{ErdosProblems.Skip.LadderT67.t83_depth_floor}{ErdosProblems/Skip/LadderT67.lean:71294}
\\
Farey denominator floor
&
Every $q\in\mathbb{N}$ with $0<q\le Q_0=79\,639\,646\,646\,701\,375\,323\,355\,774\,875\,831\,053$
satisfies the window-$(N,K)=(1,240)$ gap certificate; $q=Q_0+1$ is the exact
first failure.
&
\scale{bounded}
&
\lean{gap_check_window_1_240_le_79639646646701375323355774875831053}{GapFareyBound.lean:176}
\\
Coprimality-probability / M\"obius-square Farey floor
&
Every $d\in\mathbb{N}$ with $0<d\le Q_1=39\,819\,823\,323\,350\,687\,661\,677\,887\,437\,915\,526$
(exactly $\lfloor Q_0/2\rfloor$) is excluded as a denominator of
$\Pr(\gcd(X,Y)=1)$ and, separately, of $T=\sum_d\mu(d)/(2^d-1)^2$.
&
\scale{bounded}
&
\lean{tsum_visible_coprime_pairs_ne_int_div_of_den_le_39819823323350687661677887437915526}{CertificateKernel.lean:18760}
\lean{tsum_moebius_div_two_pow_sub_one_sq_ne_int_div_of_den_le_39819823323350687661677887437915526}{CertificateKernel.lean:18675}
\\
\end{longtable}
\normalsize

\begin{rem}[What this table is not]
No row above is cofinal in its indexing parameter ($h$, $t$, or the
denominator bound), and no row is claimed to extrapolate. The historical
diagonal-pincer depths are irregular and nonmonotone
($6,5,7,7,9,14,15,14,21,22,23,26,\dots$); no closed-form growth rate for
$D(t)$ is proved or conjectured in the corpus. This table records the
historical 28-scale bank through $t=64$; the later aggregate theorem for every
$t\le82$ is stated above and is not itemised row by row here.  Neither bounded
record supplies the cofinal obligation developed in Part~1.
\end{rem}
